\section{Context-Guided Diffusion Models}
\label{sec:context_guided_diffusion}
Context-guided diffusion models extend the classical diffusion framework by
conditioning the generative process on structured, domain-specific context.
Whereas standard conditional diffusion injects labels or text embeddings as
static side information, context-guided methods modify the \emph{trajectory} of
the reverse SDE or ODE itself. Formally, this is achieved by augmenting the
reverse dynamics
\[
\mathrm{d}x_t = 
\Big(f_t(x_t) - g_t^2 \nabla_x \log p_t(x_t)\Big)\mathrm{d}t + g_t\,\mathrm{d}\bar{w}_t
\]
with an additional context-dependent potential that steers the sampler:
\[
\mathrm{d}x_t
=
\Big(
f_t(x_t)
- g_t^2 \nabla_x \log p_t(x_t)
- \lambda \nabla_x U_{\mathrm{ctx}}(x_t \mid c)
\Big)\mathrm{d}t
+ g_t\, \mathrm{d}\bar{w}_t,
\]
where $c$ denotes contextual information and $U_{\mathrm{ctx}}$ is a potential
encoding domain-specific structure. This modification implicitly defines a new
target distribution
\[
p_{\mathrm{ctx}}(x) \propto p(x)\,e^{-U_{\mathrm{ctx}}(x\mid c)},
\]
illustrating how context restricts the generative manifold and suppresses
out-of-distribution (OOD) drift. Such mechanisms are essential in domains where
generation depends on relational, structural, or temporal constraints that
cannot be captured by prompt-based conditioning alone, including molecular
design, protein engineering, structured image generation, and video synthesis.

Across recent formulations, context-guided diffusion introduces three key
principles. First, context acts as an auxiliary latent that modifies the score
network itself. If $s_\theta(x_t,t)$ denotes the base score estimator, context
often produces an effective score of the form
\[
s_\theta(x_t,t,c)
=
s_\theta(x_t,t)
+
A_\theta(x_t,c),
\]
where $A_\theta$ is a cross-attention term encoding structural or semantic
information from the context. Second, context can be \emph{dynamic}: updated at
each sampling step through a rule $c_{t+\Delta t} = \Phi(c_t, x_t)$, allowing the
trajectory to adapt to evolving constraints. Third, context acts as a stabiliser
that improves robustness and prevents hallucinations by anchoring the sampler to
regions consistent with domain knowledge.

\paragraph{Context-Guided Diffusion for OOD Molecular and Protein Design.}
The uploaded work on molecular and protein generation is a representative example
of context as a domain-specific potential. Biochemical context---including
binding-site geometry, interaction graphs, sequence scaffolds, or structural
motifs---is encoded into a potential $U_{\mathrm{ctx}}$ that penalises steps
leading away from chemically valid configurations. By injecting the gradient
$-\nabla_x U_{\mathrm{ctx}}$ into the reverse SDE, the sampler is encouraged to
remain on the feasible biochemical manifold. This improves OOD robustness by  
(i) incorporating physically meaningful constraints at each step,
(ii) preventing drift into implausible structural regimes, and  
(iii) ensuring that generated molecules and proteins satisfy geometric and
functional requirements. These ideas translate naturally to medical imaging,
where structured radiological signals (severity scores, lesion distributions,
opacity maps) can serve as analogous forms of context.

\paragraph{In-Context Aware Image Generation (Context Diffusion).}
Context Diffusion generalises in-context learning to diffusion models by
allowing a model to ingest a ``context set''---a small number of reference
examples---and generate new samples consistent with that set. This is achieved
by embedding the context set into a latent representation and injecting it into
the score network via a cross-attention mechanism $A_\theta(x_t,c)$. The
resulting context-modulated score adapts the generative trajectory to match the
geometry and style of the context examples without additional training. This is
highly relevant for compositional diffusion, where hybrid configurations
(e.g., mixed pathologies) may lie outside all training distributions: in-context
modulation can stabilise sampling by aligning the trajectory with examples
exhibiting similar spatial patterns or semantic structure.

\paragraph{Contextualized Diffusion for Text-Guided Image and Video Generation.}
In contextualized diffusion models for text-guided image or video generation,
context is encoded hierarchically or temporally and integrated into each
denoising step through contextual cross-attention. Rather than serving as a
static conditioning token, the context modulates the drift throughout the
trajectory, enforcing global coherence across spatial layouts or temporal
sequences. Two principles emerge from this framework: (i) context must be
injected continuously along the diffusion path to maintain alignment with
high-level semantics, and (ii) its influence must be carefully calibrated so as
not to overpower the underlying generative model. These principles map directly
onto our Phase~2 experiments, where severity or lesion-location guidance must
influence the trajectory strongly enough to prevent hallucination, but not so
strongly that anatomical structure becomes distorted.

\paragraph{Relevance to Compositional Diffusion.}
Context-guided diffusion provides a crucial complement to compositional
operators. Whereas compositional methods combine distributions or score fields,
they do not inherently constrain the sampler to remain semantically consistent.
Context guidance mitigates this limitation by  
(i) stabilising the generative trajectory in regions where component models
disagree,  
(ii) imposing structured semantic constraints, and  
(iii) attenuating OOD drift through context-conditioned potentials or
cross-attention terms.  
These properties are essential for medical compositions such as TB $\land$
silicosis, where continuous radiological attributes (severity, localisation,
opacity patterns) provide natural context signals that can refine, correct, or
disambiguate the outputs of base compositional operators.

Together, these works establish context-guided diffusion as a mathematically
grounded and practically effective mechanism for integrating structured signals
into the generative trajectory. This makes it a natural foundation for our
Phase~2 investigation of property-guided compositional diffusion.
